\section{不定积分的概念}\label{sec:6.1}

微分学的基本问题是：已知一个函数，要求它的变化率，也就是求微商问题．例如，已知作变速直线运动的质点的运动规律为 $s = s(t)$，要求质点在时刻 $t$ 的瞬时速度 $v(t)$，我们只需将函数 $s = s(t)$ 对 $t$ 求微商即可，即得到 $v(t) = s'(t)$．但是在力学、物理等自然科学和许多实际问题中，往往提出相反的问题：已知作直线运动的质点在任一时刻 $t$ 的瞬时速度 $v = v(t)$，要求质点的运动规律．这种已知一个函数的微商，反回去求原来的函数的问题，就是求微商的反问题．

\begin{definition}\label{def:6-1}
设在区间 $I$ 内每一点，都有 $F'(x) = f(x)$，则称 $F(x)$ 是 $f(x)$ 在 $I$ 上的一个\emph{原函数}．
\end{definition}

\begin{example}
设 $f(x)=x^2$，则 $F(x)=\frac{1}{3}x^3$ 是 $f(x)$ 的一个原函数．显然 $\frac{1}{3}x^3+1$，$\frac{1}{3}x^3-\pi$ 也是 $f(x)$ 的原函数．更加一般地，对于任意常数 $C$，$\frac{1}{3}x^3+C$ 也是 $f(x)=x^2$ 的原函数．然而，$f(x)=x^2$ 的任意原函数是否都具有这样的形式？回答是肯定的．
\end{example}

\begin{theorem}\label{thm:6-1}
若 $F(x)$ 是 $f(x)$ 在区间 $I$ 内的一个原函数，则 $F(x) + C$ 是 $f(x)$ 的全体原函数，其中 $C$ 是任意常数．
\end{theorem}

\begin{proof}
首先，对任意给定的常数 $C$，$F(x) + C$ 是 $f(x)$ 的一个原函数，因为
\[
(F(x) + C)' = F'(x) = f(x).
\]
其次，设 $G(x)$ 是 $f(x)$ 在区间 $I$ 内的任意一个原函数，则
\[
(G(x) - F(x))' = G'(x) - F'(x) = f(x) - f(x) = 0.
\]
由微分中值定理的推论（推论~\ref{cor:5-2}），在区间 $I$ 上恒有
\[
G(x) - F(x) = C,\quad \text{即 } G(x) = F(x) + C.
\]
定理 6.1 证毕．
\end{proof}

\begin{definition}\label{def:6-2}
函数 $f(x)$ 在区间 $I$ 上的全体原函数统称为 $f(x)$ 在 $I$ 上的\emph{不定积分}，记作
\[
\int f(x)\,\mathrm{d}x.
\]
其中符号 $\int$ 称为\emph{积分号}，$f(x)$ 称为\emph{被积函数}，$f(x)\,\mathrm{d}x$ 称为\emph{被积表达式}，$x$ 称为\emph{积分变量}．
\end{definition}

若 $F(x)$ 是 $f(x)$ 的一个原函数，由定理 6.1，有
\[
\int f(x)\,\mathrm{d}x = F(x) + C,
\]
其中 $C$ 称为\emph{积分常数}．

在几何上，如果 $y = F(x)$ 是 $f(x)$ 的一条积分曲线，则对应于不同常数 $C$ 的曲线族 $y = F(x) + C$ 是由 $y = F(x)$ 沿纵轴平行移动而得到的．在相同横坐标 $x$ 处，各条积分曲线上对应点的切线斜率都等于 $f(x)$，即它们具有相同的切线斜率（见图~\ref{fig:6-1}）．

\begin{figure}[htbp]
\centering
\includegraphics[width=0.45\textwidth]{fig6-1.png}
\figcap{图 6-1}\label{fig:6-1}
\end{figure}

根据定义，求微分与求不定积分是互逆的运算：
\begin{enumerate}[(1)]
\item $\dfrac{\mathrm{d}}{\mathrm{d}x}\left[\int f(x)\,\mathrm{d}x\right] = f(x)$，或 $\mathrm{d}\left[\int f(x)\,\mathrm{d}x\right] = f(x)\,\mathrm{d}x$；
\item $\int F'(x)\,\mathrm{d}x = F(x) + C$，或 $\int \mathrm{d}F(x) = F(x) + C$．
\end{enumerate}

由微分运算法则，直接可得不定积分的线性性质：
\begin{enumerate}[(1)]
\item $\int [f(x) \pm g(x)]\,\mathrm{d}x = \int f(x)\,\mathrm{d}x \pm \int g(x)\,\mathrm{d}x$；
\item $\int k f(x)\,\mathrm{d}x = k \int f(x)\,\mathrm{d}x$（$k \neq 0$ 为常数）．
\end{enumerate}

由基本初等函数的导数公式，反过来立即得到下列\textbf{基本不定积分表}：
\begin{enumerate}[(1)]
\item $\int 0\,\mathrm{d}x = C$；
\item $\int x^\mu\,\mathrm{d}x = \dfrac{x^{\mu+1}}{\mu+1} + C\quad (\mu \neq -1)$；
\item $\int \dfrac{\mathrm{d}x}{x} = \ln|x| + C\quad (x \neq 0)$；
\item $\int a^x\,\mathrm{d}x = \dfrac{a^x}{\ln a} + C\quad (a > 0, a \neq 1)$；特别地，$\int e^x\,\mathrm{d}x = e^x + C$；
\item $\int \sin x\,\mathrm{d}x = -\cos x + C$；
\item $\int \cos x\,\mathrm{d}x = \sin x + C$；
\item $\int \sec^2 x\,\mathrm{d}x = \int \dfrac{\mathrm{d}x}{\cos^2 x} = \tan x + C$；
\item $\int \csc^2 x\,\mathrm{d}x = \int \dfrac{\mathrm{d}x}{\sin^2 x} = -\cot x + C$；
\item $\int \sec x\tan x\,\mathrm{d}x = \sec x + C$；
\item $\int \csc x\cot x\,\mathrm{d}x = -\csc x + C$；
\item $\int \dfrac{\mathrm{d}x}{1+x^2} = \arctan x + C = -\operatorname{arccot} x + C_1$；
\item $\int \dfrac{\mathrm{d}x}{\sqrt{1-x^2}} = \arcsin x + C = -\arccos x + C_1$．
\end{enumerate}

\begin{example}
求 $\int \sin^2\frac{x}{2}\,\mathrm{d}x$．
\end{example}

\begin{proof}[解]
利用半角公式 $\sin^2\frac{x}{2} = \frac{1 - \cos x}{2}$，
\[
\int \sin^2\frac{x}{2}\,\mathrm{d}x = \int \frac{1 - \cos x}{2}\,\mathrm{d}x = \frac{1}{2}\int \mathrm{d}x - \frac{1}{2}\int \cos x\,\mathrm{d}x = \frac{1}{2}x - \frac{1}{2}\sin x + C.
\]
\end{proof}

\begin{example}
求 $\int \frac{\mathrm{d}x}{\sin^2 x\cos^2 x}$．
\end{example}

\begin{proof}[解]
利用恒等式 $1 = \sin^2 x + \cos^2 x$，
\[
\int \frac{\mathrm{d}x}{\sin^2 x\cos^2 x} = \int \frac{\sin^2 x + \cos^2 x}{\sin^2 x\cos^2 x}\,\mathrm{d}x = \int \frac{\mathrm{d}x}{\cos^2 x} + \int \frac{\mathrm{d}x}{\sin^2 x} = \tan x - \cot x + C.
\]
\end{proof}

\begin{example}
求 $\int \frac{x^3 + x + 1}{1 + x^2}\,\mathrm{d}x$．
\end{example}

\begin{proof}[解]
将被积函数拆开：
\[
\frac{x^3 + x + 1}{1 + x^2} = \frac{x(x^2 + 1) + 1}{1 + x^2} = x + \frac{1}{1 + x^2}.
\]
于是
\[
\int \frac{x^3 + x + 1}{1 + x^2}\,\mathrm{d}x = \int x\,\mathrm{d}x + \int \frac{\mathrm{d}x}{1 + x^2} = \frac{1}{2}x^2 + \arctan x + C.
\]
\end{proof}

\begin{example}
求 $\int (2^x + 3^x)^2\,\mathrm{d}x$．
\end{example}

\begin{proof}[解]
展开被积表达式：
\[
(2^x + 3^x)^2 = 4^x + 2\cdot 6^x + 9^x.
\]
于是
\[
\int (2^x + 3^x)^2\,\mathrm{d}x = \int 4^x\,\mathrm{d}x + 2\int 6^x\,\mathrm{d}x + \int 9^x\,\mathrm{d}x = \frac{4^x}{\ln 4} + \frac{2\cdot 6^x}{\ln 6} + \frac{9^x}{\ln 9} + C.
\]
\end{proof}

\begin{xiti}
\item 求下列不定积分：
\begin{ziti}
\item $\int \left(x^5 + x^3 - \frac{\sqrt{x}}{4}\right)\mathrm{d}x$；
\item $\int (5 - x)^3\,\mathrm{d}x$；
\item $\int \left(\sqrt{x} + \sqrt[3]{x} + \frac{3}{\sqrt{x}} + \frac{2}{\sqrt[3]{x}}\right)\mathrm{d}x$；
\item $\int \frac{\mathrm{d}x}{x^4(1+x^2)}$；
\item $\int \frac{3x^2}{1 + x^2}\,\mathrm{d}x$；
\item $\int \frac{x + 1}{\sqrt{x}}\,\mathrm{d}x$；
\item $\int (2\sin x - 4\cos x)\mathrm{d}x$；
\item $\int (3 - \sec^2 x)\,\mathrm{d}x$；
\item $\int (\tan^2 x + 3)\mathrm{d}x$；
\item $\int \frac{2 + \sin^2 x}{\cos^2 x}\,\mathrm{d}x$；
\item $\int \frac{\tan x\,\mathrm{d}x}{\cos^2\frac{x}{2} - \sin^2\frac{x}{2}}$；
\item $\int \frac{\cos 2x}{\cos x - \sin x}\,\mathrm{d}x$；
\item $\int \frac{\mathrm{d}x}{1 + \cos 2x}$；
\item $\int (5^x + 1)^2\,\mathrm{d}x$；
\item $\int \left(2^x + \left(\frac{1}{3}\right)^x - \frac{e^x}{5}\right)\mathrm{d}x$；
\item $\int e^x\left(1 - \frac{e^{-x}}{\sqrt{x}}\right)\mathrm{d}x$；
\item $\int \left(\cos x - \frac{2}{1 + x^2} - \frac{1}{4\sqrt{1 - x^2}}\right)\mathrm{d}x$；
\item $\int \sqrt{x\sqrt{x}}\,\mathrm{d}x$；
\item $\int 2^{2x} 3^x\,\mathrm{d}x$；
\item $\int \left(\frac{3}{\sqrt{4 - 4x^2}} + \sin x\right)\mathrm{d}x$．
\end{ziti}
\item 求一曲线 $y = f(x)$，它在点 $(x, f(x))$ 处的切线斜率为 $2x$，且通过点 $(2, 5)$．
\item 已知 $f(x)$ 满足给定的关系式，试求 $f(x)$：
\begin{ziti}
\item $x f'(x) = 1\quad (x > 0)$；
\item $\frac{f'(x)}{x} = 1\quad (x > 0)$；
\item $f(x)f'(x) = 1\quad (x > 0)$；
\item $\frac{f'(x)}{f(x)} = 1\quad (f(x) > 0)$．
\end{ziti}
\end{xiti}
